\chapter{Eröffnungen}

\begin{bids}
	1\c\rs  & 11\rplus, 2\rplus \c \\
	        & \subbids{
	        a. & 12-14 balanced, ohne 5er \M \\
	        b. & 11\rplus, Einfärber/Zweifärber, längste Farbe \c \\
		    c. & 12\rplus, 4414} \\
	1\d\rs  & 11\rplus, 2\rplus \d \\
		    & \subbids{
		    a. & 18-20 balanced, ohne 5er \M \\
		    b. & 11\rplus, Einfärber/Zweifärber, längste Farbe \d \\
		    c. & 12\rplus, 4441, 4144, 1444 \\
		    d. & 11-15, 5\c 4\d, außer wenn \c s zu schön / \d s zu schwach} \\
	1\h     & 11\rplus, 5\rplus \h \\
	        & \subbids{
		    a. & 12-14 / 18-20 balanced, 5\h \\
		    b. & 11\rplus, Einfärber/Zweifärber, längste Farbe \h} \\
	1\s     & 11\rplus, 5\rplus \s\\
		    & \subbids{
		    a. & 12-14 / 18-20 balanced, 5\s \\
		    b. & 11\rplus, Einfärber/Zweifärber, längste Farbe \s} \\
	1SA     & 15-17 balanced, 5er \M{} möglich \\
	2\c\rs  & a. gf balanced und unbalanced Hände \\
			& b. 23-24	balanced, 5er \M{} möglich \\
    2\d\ra  & schwach (\tbd{in 1./2. Hand 8-10}), \M-Zweifärber \\
    2\M\rs  & 6\rplus\M, schwach \\
    2SA\rka & 21-22 balanced, 5er \M{} möglich \\
    3x      & 7\rplus x, schwach
\end{bids}

\chapter{1\c- Eröffnung}

Hände, die man mit 1\c{} eröffnet, fallen grob in zwei Kategorien: Balanced oder Unbalanced.

Alle Balanced Hände mit 12-14 Punkten und ohne 5er-Oberfarbe werden mit 1\c{} eröffnet, unabhängig davon, welche Unterfarbe länger ist.
Wenn man 1\c{} eröffnet und im Rebid eine Balanced Hand zeigt, hat man immer 12-14 Punkte,
andere Punktespannen für Balanced Hände werden stattdessen 1\d, 1SA, 2SA oder 2\c{} eröffnet.

Unbalanced Hände, die mit 1\c{} eröffnet werden sind entweder \c-Einfärber, Zweifärber mit \c{} als längste Farbe, oder 4414-Shape.

\section{Antworten auf 1\c\rs}

Auf 1er-Stufe sind die Antworten Transfers: 1\d\ra{} zeigt 4\rplus\h, 1\h\ra{} zeigt 4\rplus\s, 1\s\ra{} verneint 4er-OF und \glqq transferiert\grqq{} auf 1SA.
Mit der Balanced Hand führt man den Transfer immer aus (auch mit Fit),
mit allen anderen Händen führt man den Transfer niemals aus, sondern beschreibt seine Hand weiter.

Die Antworten 2\c{} und 2\d{} sind künstlich. 2\c\ra{} ist ein gameforcing Relay, 2\d\ra{} zeigt 5\rplus\s4\rplus\h.

Prioritäten bei den Antworten:
\begin{itemize}
	\item Mit gf Händen ohne 5er-OF reizt man 2\c\ra{} (4er-OF zeigt man später)
	\item Mit 44 in den Oberfarben ohne gf reizt man 1\d\ra
	\item Mit 5\rplus\s4\rplus\h{} antwortet man ohne gf 2\d\ra, mit gf 1\h\ra
\end{itemize}

\tbd{Alternative: Bei gf Hände mit 4er-Oberfarben diese immer auf 1er-Stufe reizen, 2\c verneint dann 4er-Oberfarbe}

Antworten auf 1\c\rs:
\nopagebreak

\begin{bids}
	Pass    & 0-5 (viele 3-5 Hände eignen sich aber für Transfers!) \\
	1\d\ra  & 3\rplus, 4\rplus\h \\
	        & mit 3-5 sollte man 5er\h{} oder min. 43 in \d\s{} haben \\
	        & mit gf und nur 4er\h{} stattdessen 2\c\ra{} reizen \\
	1\h\ra  & 3\rplus, 4\rplus\s \\
			& mit 3-5 sollte man 5er\s{} oder min. 43 in \d\h{} haben \\
			& mit gf und nur 4er\s{} stattdessen 2\c\ra{} reizen \\
	1\s\ra  & a. 6-7 balanced, keine 4\rplus\M \\
			& b. 6-11 unbalanced, keine 4\rplus OF (Notgebot) \\
	1SA     & 8-11 balanced, keine 4\rplus\M \\
	2\c\ra  & gf Relay, keine 5\rplus\M \\
	2\d\ra  & 3-11, 5\rplus\s4\rplus\h \\
	2\M     & schwach, 6\rplus\M \\
	2SA     & 11-12, 4333 ohne 4er OF \\
			& \subbids{
		      3\c & to play \\
			  Rest & gf} \\
	3\m\ra  & 9-11, 6\rplus\m{} (mind. KB10xxx), ohne 4er OF \\
			& alles außer pass ist gf \\
	3\M     & schwach, 7\rplus\M \\
	3SA     & 13-15, 4333 ohne 4er \M
\end{bids}

\subsection{1\c\rs1\d\ra}

Transfers in eine Oberfarbe führt der 1\c-Eröffner mit der Balanced Hand immer aus. Daher stehen im Rebid 1SA und 2SA als künstliche Gebote zur Verfügung.
1SA\ra{} zeigt einen Einfärber (entweder schwach mit 3er-Anschluss oder stark), 2SA\ra{} ist eine starke Hebung der Oberfarbe. Dies ermöglicht, dass 3\c{} und 4\M{} jetzt immer extreme Verteilungen zeigen und 3SA als to play zur Verfügung steht.

\begin{bids}
	Pass    & verboten \\
	1\h\ra  & 12-14 balanced, 2-4\h \\
	1\s     & 11\rplus, Zweifärber 5\rplus\c4\rplus\s, fast nie zu passen \tbd{Alternative 11-18}\\
	1SA\ra  & a. 11-15, Einfärber 6\rplus\c{} mit 3er\h \\
	        & b. 16\rplus, starker Einfärber mit 6\rplus\c \\
	2\c     & a. 11-15, Zweifärber 5\rplus\c4\rplus\d \\
	        & b. 11-15, Einfärber 6\rplus\c{} ohne 3er\h \\
	2\d     & 16\rplus, starker Zweifärber 5\rplus\c4\rplus\d, fast nie zu passen \\
	2\h     & 11-15, 3\rplus\h{} (3er fit kann nur 1345 sein, meist 4er) \\
	2\s\rka & 10\rplus, 6\rplus\c5\rplus\s \tbd{Alternative 19\rplus, 5\rplus\c4\rplus\s}\\
	2SA\ra  & 19\rplus, 4\rplus\h\\
	3\c\rka & 13-15, 7\rplus\ \\
	3\d\ra  & 19\rplus, 4\rplus\h, \d-Splinter \\
	3\s\ra  & 19\rplus, 4\rplus\h, \s-Splinter \\
	3\h     & 16-18, 4\rplus\h \\
	3SA     & to play, basierend auf semi-forcing in \c{} mit \h-Kürze \\
	4\c     & Schlemminteresse, 6\rplus\c4\h \\
	4\h     & 10-12, 6\rplus\c5\h
\end{bids}

\subsubsection{1\c\rs1\d\ra1\h\ra}

Eröffner hat 12-14 balanced und kann 2-4 \h s haben. Jetzt ist 2-Way-Checkback an.
2\c{} ist ein Relay, entweder um 2\d{} (bzw. 2\h, falls Eröffner 4er\h{} hat) zu spielen, oder um diverse einladende Hände zu zeigen.
2\d{} ist gameforcing und zeigt 5er\h{} (\tbd{nur 4\rplus\h, falls 1\c{} 2\c{} 4er-Oberfarbe verneint}).
\tbd{(Alternative: 2\c{} ist inv+ mit 5er\h, 2\d{} ist gf mit nur 4er\h)}

\begin{bids}
	Pass	& erlaubt (Plicht mit 5\rminus{} Punkten) \\
	1\s     & nicht einladend, 5\rplus\h 4\s{} \ensuremath{\rightarrow} natürlich \\
	1SA		& to play \ensuremath{\rightarrow} 2\h{} hatte 4er\h \\
	2\c\ra  & maximal inv Relay \\
	2\d\ra  & gf, 5er\h \\
	2\h     & 5er\h{}, Mini-Einladung (nur gegenüber 4er\h{} einladend) \\
	2\s     & inv, 4\h4\s \\
	2SA\ra  & Relay zu 3\c{} \ensuremath{\rightarrow} \\
			& \subbids{
			pass & wollte 3\c{} spielen \\
			Rest & SI-\c} \\
	3\m     & INV 5\h5\m{} \ensuremath{\rightarrow} pass verneint 3er\h \\
	3\h     & gf, 6er\h \\
	3\s\ra  & Splinter, 6er\h \\
	3SA		& to play \\
	4\m		& Splinter, 6er\h \\
	4\h		& to play
\end{bids}

\subsubsection{1\c\rs1\d\ra1\h\ra2\c\ra}

\begin{bids}
	2\d\ra	& kein 4er\h \\
			& \subbids{
			Pass	& will jetzt 2\d{} spielen \\
			2\h		& inv, 5er\h \\
			2\s		& inv, 5er\h4er\s \\
			2SA		& inv, 4\rplus\h \\
			3\m		& inv, 4er\h5\rplus\m \\
			3\h		& inv, 6er\h} \\
	2\h\ra	& 4er\h
\end{bids}

\subsection{1\c\rs1\h\ra}

\begin{bids}
	Pass    & verboten \\
	1\s\ra  & 12-14 balanced, 2-4\s \\
	1SA\ra  & a. 11-15, Einfärber 6\rplus\c{} mit 3er\s \\
			& b. 16\rplus, starker Einfärber mit 6\rplus\c \\
	2\c     & a. 11-15, Zweifärber 5\rplus\c4\rplus\d/\h \\
        	& b. 11-15, Einfärber 6\rplus\c {} ohne 3er\s \\
	2\d     & 16\rplus, starker Zweifärber 5\rplus\c4\rplus\d, fast nie zu passen \\
	2\h     & 16\rplus, starker Zweifärber 5\rplus\c4\rplus\h, fast nie zu passen \\
	2\s     & 11-15, 3\rplus\s{} (3er-Fit nur bei 5431-Hand-Typ, meist 4er) \\
	2SA\ra  & 19\rplus, 4\rplus\s\\
	3\c\rka & 13-15, 7\rplus\c \\
	3\d\ra  & 19\rplus, 4\rplus\s, \d-Splinter \\
	3\h\ra  & 19\rplus, 4\rplus\s, \h-Splinter \\
	3\s     & 16-18, 4\rplus\s \\
	3SA     & to play, basierend auf semi-forcing in \c{} mit \s-Kürze \\
	4\c     & Schlemminteresse, 6\rplus\c4\s \\
	4\s     & 10-12, 6\rplus\c5\s
\end{bids}

\subsubsection{1\c\rs1\h\ra1\s\ra}

\begin{bids}
	Pass	& erlaubt (Plicht mit 5\rminus{} Punkten) \\
	1SA		& to play \ensuremath{\rightarrow} 2\s{} hatte 4er\s \\
	2\c\ra  & maximal inv Relay \\
	2\d\ra  & gf, 5er\s \\
	2\h\ra  & gf, 5er\s5er\h \\
	2\s     & 5er\s{}, Mini-Einladung (nur gegenüber 4er\s{} einladend) \\
	2SA\ra  & Relay zu 3\c{} \ensuremath{\rightarrow} \\
        	& \subbids{
	    	pass & wollte 3\c{} spielen \\
	    	Rest & SI-\c} \\
	3\m     & INV 5\s5\m{} \ensuremath{\rightarrow} pass verneint 3er\s \\
	3\h\ra  & Splinter, 6er\s \\
	3\s     & gf, 6er\s \\
	3SA		& to play \\
	4\m		& Splinter, 6er\s \\
	4\s		& to play
\end{bids}

\subsubsection{1\c\rs1\h\ra1\s\ra2\c\ra}

\begin{bids}
	2\d\ra	& kein 4er\h \\
	& \subbids{
		Pass	& will jetzt 2\d{} spielen \\
		2\h		& inv, 5er\h \\
		2\s		& inv, 5er\h4er\s \\
		2SA		& inv, 4\rplus\h \\
		3\m		& inv, 4er\h5\rplus\m \\
		3\h		& inv, 6er\h} \\
	2\h\ra	& 4er\h
\end{bids}

\subsection{1\c\rs1\s\ra}

\begin{bids}
	Pass    & verboten \\
	1SA     & 12-14 balanced \\
	2\c     & a. 11-15, Zweifärber 5\rplus\c4\rplus\d/\h/\s \\
	        & b. 11-15, Einfärber 6\rplus\c \\
	2\d     & 16\rplus, starker Zweifärber 5\rplus\c4\rplus\d, forcing \\
	2\h     & 16\rplus, starker Zweifärber 5\rplus\c4\rplus\h, forcing \\
	2\s     & 16\rplus, starker Zweifärber 5\rplus\c4\rplus\s, forcing \\
	2SA\ra  & 16\rplus, starker Einfärber mit 6\rplus\c \\
	3\c     & 13-15, 7\rplus\c \\
	3x      & gf, 6\rplus\c5\rplus x
\end{bids}

\subsection{1\c\rs1SA}

\begin{bids}
	Pass    & erlaubt (Pflicht mit 12-14 bal) \\
	2\c     & a. 11-15, Zweifärber 5\rplus\c4\rplus\d/\h/\s \\
	        & b. 11-15, Einfärber 6\rplus\c \\
	2\d     & 16\rplus, starker Zweifärber 5\rplus\c4\rplus\d, gameforcing \\
	2\h     & 16\rplus, starker Zweifärber 5\rplus\c4\rplus\h, gameforcing \\
	2\s     & 16\rplus, starker Zweifärber 5\rplus\c4\rplus\s, gameforcing \\
	2SA\ra  & 16\rplus, starker Einfärber mit 6\rplus\c, gameforcing  \\
	3\c     & 13-15, 7\rplus\c \\
	3x      & 13-15 6\rplus\c5\rplus x \\
	3SA     & to play
\end{bids}

\subsection{1\c\rs2\c\ra}

\begin{bids}
	Pass    & verboten \\
	2\d\ra  & 12-14 balanced oder 4414 \\
	2\h     & a. 11\rplus, 5\rplus\c4\rplus\h \\
        	& b. 15\rplus, 4414 \\
	2\s     & 11\rplus, 5\rplus\c4\rplus\s \\
	2SA\ra  & 11\rplus, 5\rplus\c4\rplus\d \\
	3\c     & 11\rplus, 6\rplus\c{} Einfärber
\end{bids}

\subsection{1\c\rs2\d\ra}

\begin{bids}
	Pass    & nicht empfohlen \\
	2\h     & 11-17, 3\rplus\h \\
	2\s     & 11-17, 2\rplus\s \\
	2SA     & 17-19 einladend, max. 2er\s{} und 3er\h{} \\
	3\c     & to play, 6\rplus\c{} \\
	3\d\ra  & 4.FF gf \\
	3\h     & 16-18, 4\rplus\h, einladend \\
	3\s     & 16-18, 3\rplus\s, einladend \\
	3SA     & to play \\
	4\M     & to play
\end{bids}

\section{1\c\rs{} mit Gegenreizung}

Nach 1\c\rs{} (1\d/\h) ist X das \glqq gestohlene Gebot\grqq, transferiert also auf \h/\s.

\chapter{1\d- Eröffnung}

Hände, die man mit 1\d{} eröffnet, fallen grob in zwei Kategorien: Balanced oder Unbalanced.

Alle Balanced Hände mit 18-20 Punkten und ohne 5er-Oberfarbe werden mit 1\d{} eröffnet, unabhängig davon, welche Unterfarbe länger ist.
Wenn man 1\d{} eröffnet und im Rebid eine Balanced Hand zeigt, hat man immer 18-20 Punkte,
andere Punktespannen für Balanced Hände werden stattdessen 1\c, 1SA, 2SA oder 2\c{} eröffnet.

Unbalanced Hände, die mit 1\d{} eröffnet werden sind entweder \d-Einfärber, Zweifärber mit \d{} als längste Farbe, oder 4441-Hände mit 4er\d.

\section{Antworten auf 1\d\rs}

Die Antworten 2\c{} und 2\d{} sind künstlich. 2\c\ra{} ist ein gameforcing Relay, 2\d\ra{} zeigt 5\rplus\s4\rplus\h.

Prioritäten bei den Antworten:
\begin{itemize}
	\item Mit gf Händen ohne 5er-OF reizt man 2\c\ra{} (4er-OF zeigt man später)
	\item Mit 44 in den Oberfarben ohne gf reizt man 1\h
	\item Mit 5\rplus\s4\rplus\h{} antwortet man ohne gf 2\d\ra, mit gf 1\s
\end{itemize}

\tbd{Alternative: Bei gf Hände mit 4er-Oberfarben diese immer auf 1er-Stufe reizen, 2\c verneint dann 4er-Oberfarbe}

Antworten auf 1\d\rs:
\nopagebreak

\begin{bids}
	Pass    & 0-5 \\
	1\h     & 3\rplus, 4\rplus\h \\
	        & mit 3-5 sollte man 5er\h{} haben \\
        	& mit gf und nur 4er\h{} stattdessen 2\c\ra{} reizen \\
	1\s     & 3\rplus, 4\rplus\s \\
	        & mit 3-5 sollte man 5er\s{} haben \\
	        & mit gf und nur 4er\s{} stattdessen 2\c\ra{} reizen \\
	1\s\ra  & a. 6-7 balanced, keine 4\rplus\M \\
        	& b. 6-11 unbalanced, keine 4\rplus OF (Notgebot) \\
	1SA     & 6-11 keine 4\rplus\M, bal o. unbal. Notgebot \\
	2\c\ra  & gf Relay, keine 5\rplus\M \\
	2\d\ra  & 3-11, 5\rplus\s4\rplus\h \\
	2\M     & schwach, 6\rplus\M \\
	2SA     & 11-12, 4333 ohne 4er OF \\
			& \subbids{
				3\d & to play \\
				Rest & gf} \\
	3\m\ra  & 9-11, 6\rplus\m{} (mind. KB10xxx), ohne 4er OF \\
			& alles außer pass ist gf \\
	3\M     & schwach, 7\rplus\M \\
	3SA     & 13-15, 4333 ohne 4er \M
\end{bids}

\subsection{1\d\rs1\h}

Nach einer 1\M{}-Antwort reizt der 1\d-Eröffner mit der Balanced Hand immer 1SA. Im Vergleich zur 1\c-Eröffnung steht im Rebid 1SA daher nicht als künstliches Gebot zur Verfügung.
Jetzt werden 2SA\ra{} und 3\d\ra{} verwendet, um starke Einfärber zu zeigen. Dies bedeutet wiederum, dass man 4\M{} und 3SA\ra{} nun als starke Hebungen der Oberfarben benötigt. 3\d\ra{} und 4\M{} sind demnach nicht auf extreme Verteilungen beschränkt und 3SA\ra{} ist nicht to play.

\begin{bids}
	Pass    & verboten \\
	1\s     & 11\rplus, 4\rplus\rplus\d4\rplus\s, fast nie zu passen, \tbd{Alternative 11-18} \\
	1SA     & 18-20 balanced, 2-4\h \\
	2\c     & a. 11-18, 5\rplus\d4\rplus\c \\
			& b. 11-15, 5\rplus\c4\rplus\d \\
	2\d     & 11-15, Einfärber 6\rplus\d \\
	2\h     & 11-15, 3\rplus\h{} (3er fit kann nur eine 5431-Hand sein, meist 4er) \\
	2\s\rka & 10\rplus, 6\rplus\d5\rplus\s, \tbd{Alternative 19\rplus, 5\rplus\d4\rplus\s} \\
	2SA\ra  & a. 19\rplus, 6\rplus\d, Einfärber \\
	        & b. 16-18, 6\rplus\d, Einfärber mit 3er\h \\
	3\c     & 19\rplus, 5\rplus\d4\rplus\c gf. \\
	3\d\rka & 16-18, 6\rplus\d, Einfärber mit max. 2er\h \\
	3\h     & 16-18, 4\rplus\h \\
	3\s\ra  & 19\rplus, 4\rplus\h, \s-Splinter \\
	4\c     & 19\rplus, 4\rplus\h, \c-Splinter \\
	3SA\ra  & 22\rplus, 4\rplus\h \\
	4\d     & Schlemminteresse, 6\rplus\d4\h \\
	4\h     & 19-21, 4\rplus\h
\end{bids}

\subsection{1\d\rs1\s}

\begin{bids}
	Pass    & verboten \\
	1SA     & 18-20 balanced, 2-4\s \\
	2\c     & a. 11-18, 5\rplus\d4\rplus\c \\
	        & b. 11-15, 5\rplus\c4\rplus\d \\
	2\d     & a. 11-15, Einfärber 6\rplus\d \\
			& b. 11-15, 5\rplus\d4\rplus\h \\
	2\h     & 16\rplus, 5\rplus\d4\rplus\h \\
	2\s     & 11-15, 3\rplus\s{} (3er fit kann nur eine 5431-Hand sein, meist 4er) \\
	2SA\ra  & a. 19\rplus, 6\rplus\d, Einfärber \\
	        & b. 16-18, 6\rplus\d, Einfärber mit 3er\s \\
	3\c     & 19\rplus, 5\rplus\d4\rplus\c gf. \\
	3\d\rka & 16-18, 6\rplus\d, Einfärber mit max. 2er\s \\
	3\s     & 16-18, 4\rplus\s \\
	3\h\ra  & 19\rplus, 4\rplus\s, \h-Splinter \\
	4\c     & 19\rplus, 4\rplus\s, \c-Splinter \\
	3SA\ra  & 22\rplus, 4\rplus\s \\
	4\d     & Schlemminteresse, 6\rplus\d4\s \\
	4\s     & 19-21, 4\rplus\s
\end{bids}

\subsection{1\d\rs1SA}

\begin{bids}
	Pass    & erlaubt insbes. mit 4441 und \c-Kürze mit 11-15 \\
	2\c\ra  & 16\rplus, Zweifärber 5\rplus\d4\rplus\c{} ((41)44 möglich?)\\
	2\d     & 11-15, 5\rplus\d ohne 4\rplus\c \\
	2\h     & 16\rplus, 5\rplus\d4\rplus\h{} (4-4-4-1 möglich) \\
	2\s     & 16\rplus, 5\rplus\d4\rplus\s \\
	2SA\ra  & 19\rplus, 6\rplus\d, gameforcing  \\
	3\c 	& 11-15, 5\rplus4\rplus\m\m{} ((41)44 möglich) \\
	3\d		& 16-18, 6\rplus\d \\
	3\M     & 6\rplus\d5\rplus\M, gf \\
	3SA     & to play
\end{bids}

\subsection{1\d\rs2\c\ra}

\begin{bids}
	Pass    & verboten \\
	2\d\ra  & 18-20 balanced \\
	2\h     & 11\rplus, 5\rplus\d4\rplus\h oder 4-4-4-1 / 1-4-4-4 \\
	2\s     & 11\rplus, 5\rplus\d4\rplus\s oder 4-1-4-4 \\
	2SA\ra  & 11\rplus, 5\rplus4\rplus\m\m \\
	3\c     & 11\rplus, 6\rplus\d{} Einfärber \\
	3\d		& 16\rplus, 6\rplus\d{} Einfärber (sehr gute Farbe)
\end{bids}

\subsection{1\d\rs2\d\ra}

\begin{bids}
	Pass    & erlaubt \\
	2\h     & 11-17, 3\rplus\h \\
	2\s     & 11-17, 2\rplus\s \\
	2SA     & 17-19 einladend, max. 2er\s{} und 3er\h{} \\
	3\c\ra  & 4.FF gf \\
	3\h     & 16-18, 4\rplus\h, einladend \\
	3\s     & 16-18, 3\rplus\s, einladend \\
	3SA     & to play \\
	4\M     & to play
\end{bids}

\chapter{1\M-Eröffnung}

\section{Antworten auf 1\M}

\section{1\M{} mit Gegenreizung}

\begin{itemize}
	\item \tbd{Transfers nach 1\M{} (X) ?}
	\item \tbd{Transfers nach 1\M{} (2\c) ?}
	\item \tbd{1\s{} (2\d) 2\h{} forcing oder non-forcing?}
\end{itemize}

\chapter{1SA-Eröffnung}

\section{Antworten auf 1SA\rs}

\begin{table}[h]
	\centering
	\label{1SA}
	\makebox[\textwidth][c]{
		\begin{tabular}{c | c c c}
			& schwach & einladend & gf+  \\[5pt]\hline\rule{0pt}{30pt}
			(3\rminus3\rminus) 	& \makecell{pass \\ Transfer auf 3m}
			& \makecell{2SA mit 5\rplus5\rplus mm \\ 2\s{} sonst}
			& \makecell{2SA mit 5\rplus5\rplus mm \\ 2\s{} geht \\ 3\c{} geht auch} \\[25pt]
			(43\rminus)	& \makecell{pass \\ 2\c{} mgl mit 5\rplus m \\ Transfer auf 3m}	
			& \makecell{2\c{} (fit: raise, kein fit: 2SA)} 
			& \makecell[c]{3SA \\ (41)44: 3M \\ 3\c{} ohne SI\\ 2\c{} \\2\s{} mit Kürze in M} \\[35pt]
			44			& \makecell{pass \\ 2\c{} mgl mit 5\rplus m \\ Transfer auf 3m} 	
			& \makecell{2\c{} (fit: raise, kein fit: 2SA)} 
			& \makecell{2\c{} (fit: other major für SI) } \\[30pt]
			(53\rminus) 	& Jacobi 
			& \makecell{Jacobi (dann Trnf mit 5m) \\ 2\c{} mit 5\s{}3\h{} oder 5\s{} ohne 5\rplus m} 	
			& \makecell{Jacobi (dann Trnf mit 4m/3SA/raise) \\ 2\s{} mit Kürze in M \\ 3SA} \\[30pt]
			(54) 		& Garbage Stayman	
			& \makecell{2\c{} (auf 2\d: 2\s{} oder 2SA\ra) \\ nicht Jacobi!} 
			& 2\c{} (auf 2\d: 3M in kurzem M) \\[25pt]
			55 		& Garbage Stayman 
			& 3\d 
			& 4\c \\[20pt]
			6\rplus	& Jacobi
			& Jacobi (rebid: 3M-1)
			& \makecell{Jacobi (rebid: 3M-1) \\ 4M-1}
		\end{tabular}
	}
	\caption{Übersichtstabelle für die Antworten auf 1SA}
\end{table}

\begin{bids}
	2\c\rs & Stayman \tbd{ab wie vielen Punkten?}\\
	2\d\rs & \h-Transfer \tbd{nicht mit inv+ 5+\h4\s} \\
	2\h\rs & \s-Transfer \tbd{nicht mit inv+ 5+\s4\h{} oder inv+ 5+\s ohne 5+m}\\
	2\s\ra & min oder max? oder \c-Transfer \\
	2SA\ra & \tbd{pick a minor oder \d-Transfer} \\
	3\c\ra & Frage nach 5er-OF \\
	3\d\ra & 55MM, inv+ \\
	3\M    & \tbd{(41)44} \\
	3SA    & to play \\
	4\c\ra & 55MM gf oder slamforcing \\
	4\d\ra & 6+\h \\
	4\h\ra & 6+\s \\	
\end{bids}


Siehe auch die Tabelle \ref{1SA}.

Was in der Tabelle als Transfer auf 3m steht ist je nach Situation 2\s{} mit Ziel 3\c{} oder 2SA pick a minor / Ziel 3\d.

Mit Schlemminteresse ist 3\c{} ist nicht so toll, da man eventuell keine Kontrollen mehr abfragen kann. Am liebsten wendet man es an mit genau partieforcierend und (34) oder (33), ansonsten ist 2\c{} angenehmer.

\newpage

\section{1SA{} 2\c{} Stayman}

Mit beiden Oberfarben zu viert reizt Eröffner 2\h. Die Antwort 2\s{} zeigt dann eine einladende Hand mit 5\s{} und 3\rminus\h.

Von Seite des Antwortenden reizen alle einladende Hände vom Typ (4x)(5y) ohne OF-fit einfach 3m. Einladende Hände mit 5\s{} geben immer 2\s. Als Antwort auf 2\d{} beinhaltet 2SA\ra{} die einladende 5\h4\s{}-Hand.

\begin{bids}
	2\d\ra & keine 4er M \\
	& \subbids{
		2\h & garbage Stayman 5+4+MM; kein invite \\
		2\s & inv 5\s{} mit oder ohne 4\h{} (aber keine 5\h)\\
		2SA\ra & inv (5\h{} 4\s{} ist möglich) $\rightarrow$ alles außer pass nimmt an \\
		& \subbids{3\h & max und 3er\h} \\
		3m 	& to play, 5+m \\
		3\h\ra & gf; 4\h{} 5+\s \\
		3\s\ra & gf; 4\s{} 5+\h  \\
		3SA & to play \\
		4m 	& gf; 44(41) cRKCB \\
		4M 	& to play mit 6er M
	} \\
	2M 	& 4+M \\
		& \subbids{
			2SA & inv mit 4M' \\
			3m	& einladend ohne fit mit 5+m \\
			3M	& inv mit 4+M \\
			3M'\ra	& SI in M (hat 4+4MM, da nicht 3\c gestartet) \\
			3SA	& to play
		}
\end{bids}


\section{1SA 2\d/2\h}

Hier zunächst für 1SA{} 2\d. Antworten der eröffnenden Hand:

\begin{bids}
	2\h & Normalgebot \\
	3\h & Minimum mit 4er\h{} (the law) \\
	2SA & Maximum mit 4er\h. Darauf ist 3\d{} retroforcing
\end{bids}

Antworten auf 1SA\rs2\d\rs2\h:

\begin{bids}
	pass 	& erlaubt \\
	2\s\ra	& invite 5\h, keine 5+m \\
	2SA\ra 	& Transfer, 5+\c{} oder SI \\
	3\c\ra 	& Transfer, 5+\d{} oder SI \\
	3\d\ra	& inv+ mit 6+\h \\
	3\h 	& gf 5332 Hand; Frage \h-Fit? Antworte positiv mit cue-bids \\
	3SA		& choice of games \\
	3\s, 4m	& Splinter, 6+\h
\end{bids}

Die Transfers äußern dabei
	\begin{itemize}
		\item invite 5\h{} 5m, 
		\item gf 5\h{} 5m, oder
		\item SI 5\h{} 4m.
	\end{itemize}
Antworten auf die Transfers:

\begin{bids}
	3\h 		& \h-fit; nun gf \\
	Ausführung  & kein \h-fit, minimum \\
				& \subbids{
					4m		& gf; 5+m \\
					1 step 	& SI 5\h{} 4+m. Eröffner kann mit 3SA ablehnen} \\
	3SA			& kein \h-fit, maximum
\end{bids}

Für 1SA{} 2\h{} ist alles genau analog. Die einladende Hand mit 5\s{} ohne 5+m hat nach dem Normalgebot kein rebid und ist deshalb in Stayman.


\section{1SA 2\s}

2\s{} fragt min oder max. Genauer kann es sein
\begin{itemize}
	\item einladend ohne 4+\M,
	\item gf mit einer Kürze in \M{} und Zweifel an 3SA, oder
	\item Transfer auf 3\c.
\end{itemize}
Die Antworten sind

\begin{bids}
	2SA\ra 	& min\\
	3\c\ra 	& max
\end{bids}

Nach 2SA ist 3\c{} to play, ansonsten alles genauso wie die Antworten auf 3\c:

\begin{bids}
	pass 	& wollte 3\c{} spielen \\
	3\d		& deckst du \h? \\
			& \subbids{
				3SA 	& jaa! \\
				3\h\ra 	& nein und 3\rminus\s \\
				3\s\ra	& nein und 4\rplus\s 
			} \\
	3\h\ra 	& deckst du \s? ohne 3er \h{} (also eine 6+m) \\
			& \subbids{
				3SA 	& jaa! \\
				3\s\ra	& nein :( $\rightarrow${} 4m oder 5m 
			} \\
	3\s\ra 	& deckst du \s? habe 3er \h \\
			& \subbids{
				3SA 	& jaa! \\
				4m 		& nein und 4+m \\
				4\h 	& nein und zähneknirschend 4+\h
			} \\
	3SA		& to play \\
	4m		& cRKCB 
\end{bids}


\section{1SA 2SA}

2SA fragt nach der Lieblingsunterfarbe des Eröffners oder möchte 3\d{} spielen. Die Antworten sind

\begin{bids}
	3\c & \c{} sind länger oder 33mm\\
	3\d & \d{} sind länger oder 44mm
\end{bids}

Darauf wird geantwortet mit

\begin{bids}
	3\d		& to play \\
	3M		& Stopperfrage (sonst spielen wir UF) \\
	4m		& cRKCB
\end{bids}


\section{1SA 3\c}
Das 3\c-Gebot fragt nach 5er \M{} und ist partie-forcing. Nicht anzuwenden mit 4\rplus4\rplus MM, da man die nicht zeigen kann.

\begin{bids}
	3\d\ra 	& keine 5er \M \\
			& \subbids{
				3\h & 4er \s \\
					& \subbids{
						3\s & fit und keine Ruine \\
						3SA & kein fit \\
						4\s & fit und Ruine}\\
				3\s	& 4er \h \\
					& \subbids{
						3SA & kein fit \\
						4m 	& max und fit, cue\\
						4\h & min und fit}\\
				3SA	& to play \\
				4m 	& cRKCB\\
				4SA & quantitativ
			} \\
	3M		& 5er M \\
			& \subbids{
				3SA	& to play\\
				4M 	& to play\\
				M'	& SI in M\\
				4SA	& quantitativ}
\end{bids}


\chapter{2\c-Eröffnung}

Die 2\c\rs-Eröffnung ist die einzige starke Eröffnung außer 2NT und partieforcing,
mit Ausnahme der 23--24 balanced-Hände, die auch enthalten sind.
Nach Gegenreizung ist man jedenfalls immer im Forcing-Pass.
Der AW reizt meistens 2\d{}, was nichts aussagt, oder 2\h{} mit 5er.

\section{Wiedergebote auf 2\d}

Die Wiedergebote zeigen 5er-Farben auf der 2er-Stufe und 6er-Farben auf der 3er-Stufe.
Ausnahme: mit 5-5 in UF hat man kein anderes Gebot. Mit 5431 und langer UF im Zweifelsfall SA.
Um die beiden Punktspannen 23--24 und 25+ für Balanced-Hände auseinander zu halten,
ist 2\h{} vom Eröffner doppeldeutig (Kokish-Konvention).

\begin{bids}[lllX]
    2\c\rs  & 2\d\ra    & Wiedergebote: \\
            &           & 2\h\ra    & Kokish: natürlich oder balanced, gf \\
            &           & 2\s       & 5\rplus\s \\
            &           & 2NT       & 23--24 \\
            &           & 3m        & 6+m (zur Not 5, insbesondere mit 5-5 in UF) \\
            &           & 3NT\rka   & Zum Spielen mit stehender 7er-Farbe \\
\end{bids}

\subsection{Weiterreizung nach 2\h{}-Wiedergebot}

Der AW kann entweder ein 2\s{}-Relais abgeben, um die Doppeldeutigkeit aufzulösen,
oder eine eigene Farbe im Transfer reizen.
Dabei eine Besonderheit: Es gibt keinen Transfer auf \h{}. 
Mit \h{}s reizt man 2\s{}, wartet auf 2SA vom EÖ und reizt dann ganz normal Transfer.
Wenn der EÖ nicht 2SA sagt, sondern selber \h{}s hat, umso besser!
Daher ist 3\d{} Transfer auf \s{}. Das lässt dem EÖ den Platz, mit 3\h{} auf seinen \h{}s zu bestehen.

\begin{bids}[lllX]
    2\c\rs  & 2\d\ra    & 2\h\ra    & Antworten: \\
            &           &           & \subbids{
                                        2\s\ra  & Relais, oft mehr oder weniger balanced, oder lange \h{}s \\
                                        2NT\ra  & Transfer \c{} \\
                                        3\c\ra  & Transfer \d \\
                                        3\d\ra! & Transfer \s \\
                                    }
\end{bids}

\textbf{Nach 2\s{}-Relais}

\begin{bids}[llllX]
    2\c\rs  & 2\d\ra    & 2\h\ra    & 2\s\ra    & Wiedergebote: \\
            &           &           &           & \subbids{
                                                    2NT & 25\rplus{} balanced \\
                                                    \multicolumn{2}{@{}l}{Alles andere zeigt \h{}s:} \\
                                                    3\c\d\s & zweite Farbe \\
                                                    3\h     & Einfärber, gute Farbe \\
                                                    3NT\ra  & Einfärber, für NT geeignet \\
                                                } \\
\end{bids}

Eine Schwierigkeit ist, dass nach 3\h{} und 3NT kein Platz mehr ist, um unterschiedlich positive Fits zu zeigen:
Man kann nur auf 4\h{} heben (negativ, max.\ 1 Stich) oder ein Kontrollgebot abgeben (leicht positiv+).
Es gibt hier keine Verwechslungsgefahr mit einem natürlichen Gebot -- wenn der AW eine eigene Farbe zu reizen hat, muss er gleich transferieren.

\begin{tbdblock}
    Man kann nach dem 2\s{}-Relais 3\c{} und 3\h{} vom EÖ vertauschen, sodass mit
    Einfärbern mehr Platz bleibt, auf Kosten der \h{}-\c{}-Zweifärber.
    Oder 3\c{} doppeldeutig mit 3\d{}-Relais.
\end{tbdblock}

\textbf{Nach Transfers}

Der EÖ führt den Transfer mit der balanced-Hand aus und reizt natürlich mit einer \h{}-Hand.
Mit Fit kann er im Sprung heben, um die Trumpffarbe festzulegen, egal ob balanced oder \h{}s.
Eine „Wiederholung” von 3\h{} legt \h{} als Trumpf fest -- eine Farbe vom AW ist darauf ein Cuebid, selbst die Wiederholung seiner Transfer-Farbe.

\begin{bids}[lllX]
    2\c\rs  & 2\d\ra    & 2\h\ra    & Transfer-Antworten und Weiterreizung: \\
            &           &           & \subbids{
                                        2NT\ra  & Transfer \c{} \\
                                                & \subbids{
                                                    3\c\ra  & 25\rplus bal. \\
                                                    3\d\s   & zweite Farbe mit \h{}s \\
                                                    3\h     & Gute \h{}, legt Trumpffarbe fest \\
                                                    3NT     & nat. mit \h{}s (kein Fit, Stopper in Restfarben) \\
                                                    4\c     & Fit \\
                                                } \\
                                        3\c\ra  & Transfer \d \\
                                                & Weiterreizung analog (4\c = zweite Farbe) \\
                                        3\d\ra! & Transfer \s \\
                                                & Weiterreizung analog: \\
                                                & \subbids{
                                                    3\h     & Gute \h{}, legt Trumpffarbe fest \\
                                                    3\s\ra  & 25\rplus bal. \\
                                                    3NT     & nat. mit \h{}s \\
                                                    4\c\d   & zweite Farbe \\
                                                    4\s     & Fit, ermutigend, nonforcing \\
                                                    4NT     & Assfrage, bestätigt \s{}
                                                } \\
                                    } \\
\end{bids}

In der letzten Zeile ist 4NT Assfrage auch ohne vorherige Fitbestätigung:
Es gibt hier wenig Grund für ein quantitatives Gebot; mit ausgeglichenen Händen kann man
ja erst den Transfer ausführen.
4\s{} als forcing zu spielen klingt jedenfalls gefährlicher, auch wenn es nicht forcierend selten
gereizt und noch seltener gepasst wird.

\subsection{Weiterreizung nach natürlichen Wiedergeboten}

\textbf{Nach 2\s{}}

Anders als nach 2\h{} ist kein Platz mehr für Transfers, 2SA wird als abwartendes Relais benötigt.
Dafür kann der EÖ auch nicht mehr die balanced-Hand haben, wodurch die Weiterreizung ziemlich analog wird.
Dazu kommt nur die Möglichkeit für den AW, die \s{}s direkt zu heben, auf 3er-Stufe, 4er-Stufe, oder mit Splinter.

\begin{bids}[lllX]
    2\c\rs  & 2\d\ra    & 2\s       & Weiterreizung: \\
            &           &           & \subbids{ 
                                        2NT\ra  & Abwartend \\
                                        3\c\d\h & 5\rplus{}er, positive Hand \\
                                                & \subbids {
                                                    3\s & legt \s{} als Trumpf fest \\
                                                    4\h{} auf 3\h & Fit, ermutigend, nonforcing \\
                                                    4NT auf 3\h & Assfrage, bestätigt \h{} \\
                                                } \\
                                        \multicolumn{2}{@{}l}{Hebungen:} \\
                                        3\s     & min.\ leicht positiv (ab 1,5 Stichen) \\
                                                & \subbids {
                                                    3NT & non-serious -- so gerade eine GF-Eröffnung
                                                } \\
                                        4\c\d\h & Splinter: 4 Trümpfe, Single, ca.\ 1 nützliche Figur (idR 1 A/K, evtl.\ noch 1D dazu) \\
                                        4\s     & negativ (max.\ 1 Stich) \\
                                    } \\
\end{bids}

\tbd{Man kann 3SA als mittlere Hebung spielen, kein SI, aber auch kein totaler Schrott.
Dann werden 3\s{} und 4\s{} eindeutiger.}

\textbf{Nach 3UF}

Eine neue Farbe vom AW kann eine lange Farbe sein, oder auch nur ein Stopper um 3SA zu spielen -- letzeres v.A.\ mit schwachen Händen, die nicht am UF-Schlemm interessiert sind.
Wenn der EÖ kein besseres Gebot hat, kann er mit 3er-Anschluss heben; der AW korrigiert auf 5 in Eröffnerfarbe, wenn es keine lange Farbe war.

\section{Andere Antworten als 2\d}

\subsection*{2\h{}-Antwort}

2\h{} nimmt dem EÖ mit unausgeglichenen Verteilungen keinen Bietraum, daher sind die Anforderungen viel geringer, als bei den anderen Antworten:
Eine brauchbare 5er-Farbe in einer positiven Hand reicht.
Trotzdem ist der AW nicht gezwungen, ein 5er-\h{} zu zeigen, wenn er der Meinung ist, nach einer 2\d{}-Antwort besser dazustehen.
Die Weiterreizung ist natürlich und folgt den gleichen Prinzipien wie nach 2\d-Antwort.
2SA vom EÖ zeigt 23\rplus{} -- die Punktspanne wird also nicht mehr genauer aufgelöst, dafür hat der AW jetzt die 3er-Stufe zur Verfügung, um eine zweite Farbe zu zeigen.

\subsection*{Höhere Antworten}

Die Antworten ab 2\s{} haben ein erhöhtes Risiko, dass der EÖ seine Farbe eine Stufe höher zeigen muss.
Sie werden deshalb viel seltener und nur mit bestimmten Handtypen abgegeben:
\begin{enumerate}
\item Gute 5-5-Hände.
\item Einfärber mit sehr guten Farben, die schon gegenüber einer Single-Topfigur
stehend sind (6er mit drei Figuren, 7er mit zwei Topfiguren).
\end{enumerate}

In jedem Fall ist das Ziel, eine Fitfarbe für einen Schlemm zu finden, die man nach
2\d-Antwort nicht gefunden hätte: Mit 5-5-Händen kann man gegenüber einem Einfärber
einen 5-3-Fit finden, hätte nach 2\d{} aber nicht genug Platz, beide Farben zu reizen.
Mit sehr guten Einfärbern kann es auch gegenüber einer Kürze die beste Trumpffarbe sein,
aber man könnte die Farbe nach 2\d{} höchstens einmal reizen.

\begin{bids}[llX]
    2\c\rs  & 2\s\ra    & 5\rplus\s-5\rplus x oder Einfärber \\
            &           & \subbids{
                            2NT\ra  & Relais \\
                                    & \subbids{
                                        3x  & 5\rplus-5\rplus \\
                                        3\s & Einfärber \textrightarrow{} Cuebids \\
                                    } \\
                        } \\
            & 2NT\ra    & 5\rplus-5\rplus UF \\
            & 3\c\d     & Einfärber \\
\end{bids}

\section{Nach Gegenreizung}

Nach Gegenreizung fehlt der Bietraum, um die SA-Punktspannen genau aufzulösen.
Die Reizung ist deshalb einfach partieforcing.

\begin{tbdblock}
    Am besten haben wir allgemeine Regeln, wann X in forcierenden Situationen take-out
    ist und wann Strafe (z.B. 1\h{} (--) 2\c{} (4\d) X = ?).
    Auf niedriger Stufe finde ich Takeout logischer, um es dem Partner leichter zu 
    machen, wenn der andere Gegner hebt. Das ergäbe dann:
\end{tbdblock}

\subsection*{Reaktion vom Antwortenden auf Gegenreizung}

\begin{bids}[lllX]
    2\c\rs  & (2\s) & Neue Farbe & positiv, ordentliche Farbe \\
            &       & X         & positiv, take-out \\
            &       & p         & kein klares Gebot \\
            &       &           & \subbids{
                                    2NT     & bal.\ (im Prinzip), 23\rplus, forcing \\
                                    X       & kein Gebot (hier im Prinzip balanced ohne ausreichenden Stopper, auf höherer Stufe schwammiger) \\
                                    3NT     & zum Spielen \\
                                    3\s\ra  & Frage nach Stopper mit stehender Farbe \\
                                } \\
\end{bids}

Ein Ausspielkontra wird wie ein 2\c-Gebot behandelt:

\begin{bids}[lllX]
    2\c\rs  & (X)   & 2\d   & natürlich \\
            &       & XX    & positiv, take-out \\
            &       & p     & nichts zu sagen \\
            &       &       & \subbids{ 
                                XX & Kein natürliches Gebot \\
                                3\c & Frage nach Stopper \\
                            } \\
\end{bids}

\subsection*{Reaktion vom Eröffner nach Reizung vom rechten Gegner}

Im Prinzip ähnlich, wobei der Eröffner daran denken sollte, dass er erst mal passen kann,
wenn er sich über Fitfarbe oder Kontrakthöhe nicht sicher ist.

\begin{bids}[lllllX]
    2\c\rs  & (--)  & 2\d\ra    & (3\c) & X     & take-out, spielbereit in Restfarben \\
            &       &           &       & p     & abwartend \\
\end{bids}

\chapter{2\d-Eröffnung}
Die 2\d-Eröffnung verspricht einen OF-Zweifärber (mindestens 54) unter Eröffnungsstärke. Da der Partner mit 3er OF zunächst nicht weiß, ob er einen Fit hat, sollte das Gebot in erster und zweiter Hand eher solide sein (ca 8-10 HP mit vernünftigen Oberfarben, oder mindestens 10 Karten in OF).

\section*{Antworten und Weiterreizung}
Alle OF-Gebote sowie 3\c{} und 3SA sind to play. Passen um \d{} zu spielen ist auch möglich.
Aus 3SA kann EÖ mit 65 Händen rauslaufen: 4\c{} (6er \h) bzw. 4\d{} (6er \s).

Weitere Antworten:

\begin{bids}
2SA & Fragt. Rebids: \\
    & \subbids{
        3\c & 4-5 \\
        3\d & 5-4 \\
        3\h & 5-5 Minimimum \\
        3\s & 5-5 Maximum \\
    } \\
4\c & Schlemminteresse in \h \\
4\d & Schlemminteresse in \s \\
\end{bids}

Bei den beiden schlemminteressierten Geboten reizt EÖ mit Minimum 4 in besagter OF, mit Maximum das Zwischengebot.

\chapter{2OF-Eröffnung}
Weak-Two. Antworten:

\begin{bids}
Neue Farbe & natürlich, forcing \\
Hebung & Abschluss \\
2SA & einladend+, fragt. Rebids: \\
& \subbids{
3OF & Minimum \\
3 sonst & Werte \\
3SA & keine Werte außerhalb \\
4F & Splinter (ggf \h für \s) \\
}\\
\end{bids}

\chapter{Schlemmreizung}

\section{Cuebids und Last Train}

\begin{itemize}
	\item Cuebids zeigen Erst- oder Zweitrundenkontrolle in der Farbe.
	\item 3SA kann ein Ersatzcuebid sein, falls 3\s{} z.B. Non-Serious ist.
	\item 4OF-1 als Antwort auf ein 4OF-2 Cuebid/Splinter ist Last-Train (Kontrolle oder Extrastärke fehlt, um die Assfrage zu stellen).
\end{itemize}

\section{3M Modul}

\subsection{Non-Serious Step 1 (wenn beide unlimitiert sind)}

Wenn in einer gameforcing oder schlemmeinladenden Reizung ein Oberfarbenfit mit 3\M{} (3M) bestätigt wurde:

\begin{itemize}
	\item Step 1 (= das erste Gebot nach 3\M) ist Non-Serious, ein direktes Cuebid ist Serious.
	\item Falls \h{} die Trumpffarbe ist, ist 3\s\ra{} Non-Serious und 3SA\ra{} Serious mit \s-Cuebid.
	\item Falls \s{} die Trumpffarbe ist, ist 3SA\ra{} Non-Serious nur Semi-Forcing (Partner darf passen, wenn er lieber 3SA spielen will).
\end{itemize}

\subsection{No Shortage (serious) Step 1 (wenn mindestens einer limitiert ist.)}

Falls ein Reizer bereits stark limitiert ist (=Punkterange max 3F), so kann er nie serious sein.
Bedeutungen der Gebote vom limitierten Reizer:

\begin{bids}
	Step1   & Waiting ohne Kürze \\
	Cuebid  & Kürze \\
	4M      & Katastrophenblatt
\end{bids}

Vom nicht limitierten Reizer (mit limitiertem Partner) ist jeder Schlemmversuch serious.
Bedeutungen der Gebote vom nicht limitierten Reizer (mit limitiertem Partner):

\begin{bids}
	Step1   & serious ohne Kürze (oder will Kürze nicht zeigen) \\
	Cuebid  & Kürze \\
	4M      & to play
\end{bids}

No Shortage Step 1 entfällt, falls derjenige keine Kürze haben kann, dann ist 3SA Spielvorschlag.

\section{RKCB 1430}

Bei gefundenem Oberfarbenfit / wenn ein Partner eine Oberfarbe als Trumpf etabliert hat, ist 4SA immer RKCB.

\begin{bids}
	Step1   & 1 oder 4 \\
	Step2   & 0 oder 3 \\
	Step3   & 2 (oder 5) ohne Trumpfdame \\
	Step4   & 2 (oder 5) mit Trumpfdame \\
	\tbd{}  & \tbd{Antworten mit Chicane}
\end{bids}

Nach RKCB fragen wir mit \hyperref[subsec:Spiral Scan]{Spiral Scan} weiter.

Bedeutung von 4SA in anderen Situationen:
\begin{itemize}
	\item Sprung in 4SA ohne etablierte Trumpffarbe ist i.d.R. quantitativ.
	\item 4SA ohne Sprung kann to play sein, wenn keine Oberfarbe als Trumpf etabliert wurde.
	\item In vielen Situationen in der kompetitiven Reizung / Gegenreizung ist 4SA pick a minor oder zeigt beide Unterfarben.
\end{itemize}

\subsection{Bedingte Assfrage}

\tbd{Wollt ihr das spielen?}

\subsection{Exclusion RKCB}

\tbd{3014 statt 1430?}

\subsection{Spiral Scan} \label{subsec:Spiral Scan}

Nach der ersten Antwort werden ggf. die Trumpf Dame und dann die einzelnen Könige erfragt und gereizt.
\tbd{Bekannte Chicanes des Partners werden dabei ausgelassen, Singles nicht. Karten, die der Antwortende nicht haben kann, werden ebenfalls ausgelassen} \\
Nach Step1/2 ist das nächstmögliche Gebot, das nicht in Trumpffarbe ist, die Frage nach Trumpfdame. Mit dem niedrigsten mögliche Gebot in der Trumpffarbe verneint man.
\tbd{Trumpfdame ab 11?-Kartenfit dazulügen?; Reihenfolge der Könige; Positivantworten}